\documentclass{article}
\usepackage{amsmath,amssymb}
\usepackage{xurl}
\usepackage[hidelinks]{hyperref}
% Shared commands for notes in docs/paper-gaps/.

\DeclareUnicodeCharacter{2080}{\ifmmode{}_{0}\else\textsubscript{0}\fi}
\DeclareUnicodeCharacter{2081}{\ifmmode{}_{1}\else\textsubscript{1}\fi}
\DeclareUnicodeCharacter{2082}{\ifmmode{}_{2}\else\textsubscript{2}\fi}
\DeclareUnicodeCharacter{2099}{\ifmmode{}_{n}\else\textsubscript{n}\fi}

\newcommand{\Lean}{\textsc{Lean}}
\ExplSyntaxOn
\NewDocumentCommand{\leanid}{m}{
  \ifmmode
    \textsf{\detokenize{#1}}
  \else
    \group_begin:
    \urlstyle{sf}
    \tl_set:Nn \l_tmpa_tl {#1}
    \tl_replace_all:Nnn \l_tmpa_tl {\_} {_}
    \exp_args:NV \path \l_tmpa_tl
    \group_end:
  \fi
}
\ExplSyntaxOff
\providecommand{\tr}{\operatorname{tr}}
\newcommand{\tnleanIssueBase}{https://github.com/LionSR/TNLean/issues/}
\newcommand{\tnleanPrBase}{https://github.com/LionSR/TNLean/pull/}
\newcommand{\tnleanDocBase}{https://sirui-lu.com/TNLean/docs/}
\newcommand{\ghissue}[1]{\href{\tnleanIssueBase#1}{GitHub issue \##1}}
\newcommand{\ghpr}[1]{\href{\tnleanPrBase#1}{PR \##1}}
\newcommand{\doclink}[1]{\href{\tnleanDocBase#1.html}{\path{#1}}}

% Dirac notation and the identity operator, as in blueprint/src/macros/common.tex.
% \providecommand keeps note-local definitions authoritative.
\usepackage{dsfont}
\providecommand{\ket}[1]{|#1\rangle}
\providecommand{\bra}[1]{\langle #1|}
\providecommand{\braket}[2]{\langle #1 | #2 \rangle}
\providecommand{\ketbra}[2]{|#1\rangle\!\langle #2|}
\providecommand{\Id}{\mathds{1}}
\providecommand{\one}{\mathds{1}}
\providecommand{\C}{\mathbb{C}}
\providecommand{\MN}[1]{M_{#1}(\C)}
\providecommand{\MD}{M_D(\C)}

% External referencing for paper sources.  The build helper writes sanitized
% auxiliary files under build/paper-aux/ and excludes duplicated source labels.
\usepackage{xr}

% Import a generated paper auxiliary file with a caller-supplied label prefix.
% The candidate paths support builds from docs/paper-gaps/, the repository root,
% and build/paper-aux/ itself.
\newcommand{\importpaperaux}[2]{%
  \IfFileExists{../../build/paper-aux/#2.aux}{%
    \externaldocument[#1]{../../build/paper-aux/#2}%
  }{%
    \IfFileExists{build/paper-aux/#2.aux}{%
      \externaldocument[#1]{build/paper-aux/#2}%
    }{%
      \IfFileExists{#2.aux}{%
        \externaldocument[#1]{#2}%
      }{}%
    }%
  }%
}

% General external reference macros
\newcommand{\paperref}[2]{\ref{#1-#2}}
\newcommand{\papereqref}[2]{(\ref{#1-#2})}

% CPSV16 source (1606.00608 / MPDO-22-12-17-2.tex) external document and shortcuts
\importpaperaux{cpsv16-}{1606.00608/MPDO-22-12-17-2-tnlean}
\newcommand{\cpsvref}[1]{\paperref{cpsv16}{#1}}
\newcommand{\cpsveqref}[1]{\papereqref{cpsv16}{#1}}

% Verdict marker: \gapnote{<kind>}{<status>}. Kind is the policy
% classification (clarification, local-correction, scope-restriction,
% unfaithful, false-source, open-gap); status is open, wip (elimination
% underway), resolved, or historical. Machine-read by the site index;
% prints a verdict line.
\newcommand{\gapnote}[2]{\par\noindent\textbf{Verdict:} #1 (#2).\par}

\title{Fibonacci String-Net: Entries of the Two Operator Blocks}
\author{}
\date{2026-09-24}
\begin{document}
\maketitle
\gapnote{local-correction}{resolved}

\section{What the source prints}
Appendix D.1.1 of \cite{Bultinck2017Anyons} gives the categorical data of the
Fibonacci theory: labels $1,\tau$ with
$N_{11}^1=N_{\tau1}^\tau=N_{1\tau}^\tau=N_{\tau\tau}^1=N_{\tau\tau}^\tau=1$,
quantum dimensions $d_1=1$, $d_\tau=\phi=(1+\sqrt5)/2$, the F-symbols
\begin{align}
    [F^{abc}_d]_e^f & =\delta_{abe}\delta_{cde}\delta_{adf}\delta_{bcf}\,F^{abc}_{def},
    \qquad
    [F^{\tau\tau\tau}_\tau]
    =\begin{pmatrix}1/\phi & 1/\sqrt\phi\\ 1/\sqrt\phi & -1/\phi\end{pmatrix},
    \label{eq:fib_f}
\end{align}
all other admissible components being $1$, and the G-symbols
$G^{abc}_{def}=F^{abc}_{def}/(v_ev_f)$ with $v_i=\sqrt{d_i}$. The operator
tensor itself is given only as a diagram. The text then states that, after the
zero rows and columns are removed, the projector matrix product operator has
bond dimension $5$ and consists of two blocks $B_1,B_\tau$ of dimensions $2$
and $3$ that satisfy the Fibonacci fusion rules, with weights
$w_1=1/(1+\phi^2)$ and $w_\tau=\phi/(1+\phi^2)$.\footnote{Source traceability:
    \path{References/1511.08090/AnyonsPEPS.tex}, lines 1240--1269.}
No numeric entry of $B_1$ or $B_\tau$ is printed.

\section{The entries used}
Write a physical letter as a pair $(x',x)$ of an outgoing and an incoming label
and let the bond index at site $j$ run over the admissible letters
$(x'_j,x_j)\in\{(1,\tau),(\tau,1),(\tau,\tau)\}$. The $\tau$ block used in the
formalization is
\begin{align}
    (B_\tau^{(x',x)})_{(x'_j,x_j),(x'_{j+1},x_{j+1})}
     & =\delta_{(x'_j,x_j),(x',x)}\,
    [F^{\tau x\tau}_{x'_{j+1}}]_{x'}^{x_{j+1}},
    \label{eq:fib_block}
\end{align}
with the bare F-symbols of (\ref{eq:fib_f}) and without the closed-loop factors
$v_a$. The trivial block $B_1$ is the admissibility projector of bond
dimension $2$, carrying the matrix $\begin{pmatrix}0&1\\1&1\end{pmatrix}$ on its
diagonal letters. The placement (\ref{eq:fib_block}) was fixed by an
exhaustive search over the assignments of the six F-symbol slots to the
available labels, keeping those for which the fusion rules hold at lengths
$2,3,4$.\footnote{Verification record:
    \path{Notes/OpenProblemsTN/checks/p5_more_examples_data.md}, \S3.1. It is
    a record of the search, not a source.}

\section{What is checked}
Every entry of the $\tau$ block equals the corresponding F-symbol of
(\ref{eq:fib_f}); the blocks have the recorded dimensions $2$ and $3$; their
periodic operators satisfy $O_1^2=O_1$, $O_1O_\tau=O_\tau O_1=O_\tau$ and
$O_\tau^2=O_1+O_\tau$ at every positive length; and
the trace of the bond-five tensor $B_1\oplus B_\tau$ against
$\Delta=w_1 I_2\oplus w_\tau I_3$ with the printed weights is
$P_L=w_1O_1+w_\tau O_\tau$, which is idempotent at every positive length.\footnote{Results:
    \leanid{FibonacciCompression.fibTauGolden_eq_fSymbol},
    \leanid{FibonacciCompression.isMPOFusionAlgebra_fibBlock},
    \leanid{FibonacciCompression.fibProjector_eq_trace_fibDelta},
    \leanid{FibonacciCompression.fibProjector_mul_self}.}
Since the source prints no numeric entries, what is recorded is a choice of
explicit blocks with the printed block dimensions, weights and fusion rules,
not an error in the source.

\section{Relation to the source's tensor}
The source draws one site of its operator as the crossing of the horizontal
operator line $f$ with a vertical edge line $e$, with corner plaquettes $b$
(upper left), $a$ (upper right), $c$ (lower left) and $d$ (lower right), and
gives it the entry
\begin{align}
    S_f\bigl[(b,e),(c,e);(b,c),(a,d)\bigr]
     & =G^{abc}_{def}\sqrt{v_av_bv_cv_d},
    \qquad G^{abc}_{def}=\frac{F^{abc}_{def}}{v_ev_f},\quad v_i=\sqrt{d_i},
    \label{eq:fib_source}
\end{align}
where a physical letter records the left plaquette and the edge label, the
edge label passes through the site unchanged, and the bond letters are the
admissible pairs of plaquettes above and below the operator line; the
upper-right plaquette of a site is the upper-left plaquette of the next
one.\footnote{Source traceability:
    \path{References/1511.08090/AnyonsPEPS.tex}, lines 1044--1050 and
    1257--1268, figure file \path{StringNetMPO_RHS.pdf}.}
The Fibonacci G-symbols have the tetrahedral symmetry
$G^{abc}_{def}=G^{fce}_{abd}$. On the configurations whose edge labels are all
$\tau$ this moves the G-symbol of (\ref{eq:fib_source}) to the F-symbol slots
of (\ref{eq:fib_block}), and the remaining factors are
$\sqrt{v_av_bv_cv_d}/(v_bv_d)=h(a,d)/h(b,c)$ with
$h(u,l)=(v_u/v_l)^{1/2}$. Hence the diagonal bond matrix $h$ conjugates every
letter of the source's tensor on edge label $\tau$ into the corresponding letter
of $B_1$ or $B_\tau$,
\begin{align}
    h\,S_f^{(x',x)}\,h^{-1} & =B_f^{(x',x)},
    \label{eq:fib_conj}
\end{align}
the entries of $S_f$ involving $\phi^{\pm1/4}$ while those of $B_f$ lie in
$\mathbb{Z}[\phi^{-1/2}]$. The similarity cancels in the trace, so on edge label
$\tau$ the periodic operators of the source's blocks equal those of $B_1$ and
$B_\tau$ at every length, and they satisfy the Fibonacci fusion
rules.\footnote{Results:
    \leanid{FibonacciCompression.fibStringNetEdgeTau_conj},
    \leanid{FibonacciCompression.mpo_fibStringNetTensor_edgeTau},
    \leanid{FibonacciCompression.isMPOFusionAlgebra_fibStringNetEdgeTau}.}
The operator statements are restricted to edge label $\tau$: when an edge
carries the label $1$ the two plaquettes on either side of it coincide, and the
fusion rules of the source's tensor on the full alphabet of four letters per
site (plaquette and edge label) were checked numerically for one to four sites
but are not proved.

The review \cite{Cirac2021Matrix} writes the operator tensor with plaquettes
$A$ (upper left), $B$ (upper right), $C$ (lower left), $D$ (lower right), and
entry $(1/\sqrt{d_Ad_D})\,(F^{aC\alpha}_B)^D_A$.\footnote{Review traceability:
    \path{Papers/2011.12127/TN-Review-main.tex}, lines 2613--2625, figure
    file \path{fig3_mpo.pdf}.}
At $\alpha=\tau$ this is the bare block (\ref{eq:fib_block}) multiplied by
$1/\sqrt{d_{x'}d_D}$, with $D$ the lower-right plaquette. The bond gauge $g(u,l)=d_l^{-1/2}$ turns this factor
into the letterwise scalar $(d_{x'}d_x)^{-1/2}$, and the periodic operator is
the congruence
\begin{align}
    O_f^{\mathrm{rev}} & =\Delta^{-1/2}\,O_f\,\Delta^{-1/2},
    \qquad \Delta=\operatorname{diag}\prod_k d_{x_k},
    \label{eq:fib_review}
\end{align}
not a similarity of $O_f$. Read in the orthonormal basis of plaquette
configurations, these operators do not satisfy the fusion rules: at one site
$O_1^{\mathrm{rev}}=\operatorname{diag}(0,1/\phi)$, whose square has the entry
$1/\phi^2\neq1/\phi$. The operators
$O_f^{\mathrm{rev}}\Delta=\Delta^{-1/2}O_f\Delta^{1/2}$ are similar to $O_f$
and do satisfy them, so the review's normalization gives the fusion algebra
with respect to the $\Delta$-weighted inner product.\footnote{Results:
    \leanid{FibonacciCompression.mpo_fibReviewTensor},
    \leanid{FibonacciCompression.not_isMPOFusionAlgebra_fibReviewTensor},
    \leanid{FibonacciCompression.isMPOFusionAlgebra_fibReviewWeightedTensor}.}
The review makes no fusion statement for its prefactored tensor at this
passage, so this is a remark on normalization, not an error in the review.

\section{The empty chain}
The source asks for $P_L^2=P_L$ at every length; the exclusion of the empty
chain $L=0$, where $P_0$ is a scalar different from $0$ and $1$, is recorded
separately in \path{docs/paper-gaps/bmwshv17_fibonacci_projector_positive_length.tex}.

\bibliographystyle{alpha}
\bibliography{references}
\end{document}
